\section*{第 \ref{ch:b3:grand-canonical} 章——巨正则系综}

\begin{solution}{exo:b3:grand-canonical:1}
(a) 由密度较大的盒子流向较稀的盒子：从高 $\mu =
k_{\text{B}}T\ln(n\lambda_T^3)$ 到低。(b) 温度 $T$ 相同，则
$\lambda_T$ 相同，故 $\mu$ 相等即 $n$ 相等。(c)
$\mu_{\text{液体}} = \mu_{\text{蒸气}}$ ——
饱和湿度的定义。(d) 干燥空气中，蒸气的 $\mu$（$n$ 较低）
处于液体之下：水沿着 $\mu$ 的“山坡”下行而蒸发，
与它是否沸腾无关。
\end{solution}

\begin{solution}{exo:b3:grand-canonical:2}
(a) $\lambda_T = \qty{1.9e-11}{m}$。(b) $n\lambda_T^3 =
\num{1.7e-7}$：$\mu = k_{\text{B}}T\ln(\num{1.7e-7}) =
\qty{-0.40}{eV}$。(c) 熵项起作用：新分子可进入
$\sim1/n\lambda_T^3$ 个态，$-TS$ 胜过任何动能代价。
(d) $n = \lambda_T^{-3} \approx \qty{1.5e32}{m^{-3}}$ ——
约为空气液态密度的数千倍；$n\lambda_T^3 \sim 1$
是量子简并的边界
（见\cref{exo:b3:grand-canonical:11}）。
\end{solution}

\begin{solution}{exo:b3:grand-canonical:3}
$\Xi = 1 + \eu^{-\beta(\epsilon-\mu)}$；$\bar n = k_{\text{B}}T\,
\partial_\mu\ln\Xi = \eu^{-\beta(\epsilon-\mu)}/(1 +
\eu^{-\beta(\epsilon-\mu)})$，即所给形式。极限情形：$1$ 位于
$\mu$ 以下很深处，玻尔兹曼尾在上方很远处，恰为 $\tfrac12$
于 $\epsilon = \mu$ 处。这正是肖特基固体的热双能级占据
——同一条曲线，换了标签。
\end{solution}

\begin{solution}{exo:b3:grand-canonical:4}
(a) 每个电子 \qty{3.0}{eV}；$\times N_{\text{A}}e$：
$\approx\qty{290}{kJ/mol}$。(b) 化学势之差
即键能——以电子伏特计：电池继承了化学的能量
尺度。(c) 没有任何化学偶对能提供 \qty{12}{eV} 的台阶（否则会
撕裂电解质）；串联叠加可安全地累加 $\mu$ 的降落。
(d) 其两极的电子化学势已经
相等：没有梯度，也没有推动力，无论还剩多少物质。
\end{solution}

\begin{solution}{exo:b3:grand-canonical:5}
(a) 公比为 $\eu^{-\beta(\epsilon-\mu)} < 1$ 的几何级数；
对其求导。(b) 当 $\mu \to \epsilon_{\min}$ 时，级数
发散：该能级可吸纳无限多的粒子——
这种宏观堆积即玻色–爱因斯坦凝聚。(c)
$\bar n = 9.51$、$0.58$、$0.052$，对应玻尔兹曼的 $0.90$、
$0.37$、$0.050$：在稀薄尾部完全相同，在
$\mu$ 附近则被大幅增强。(d) 由于器壁可随意产生和吸收光子，
平衡态不能依赖于光子的数目：增加一个光子的
自由能代价为零，故 $\mu = 0$。
\end{solution}

\begin{solution}{exo:b3:grand-canonical:6}
(a) $P = P_0$；$P = 9P_0$。(b) 一个独立位点给出
双曲线 $P/(P + P_0)$——肌红蛋白；血红蛋白的四个
相互作用位点协同结合，使曲线陡化为
S 形，在肺部充分结合，\emph{并}在
组织中释放。(c) $\eu^{0.13/0.0267} \approx 130$：在分压相等时，
CO 争夺位点的能力高出百倍——微量 CO
即可驱逐血液中的氧。(d) 提高 O$_2$ 分压
会提高其 $\mu$，并改变竞争的走向：等温线
争夺的是 $\mu$，而高压氧只是
“出价”更高。
\end{solution}

\begin{solution}{exo:b3:grand-canonical:7}
(a) 对 $\langle N\rangle = k_{\text{B}}T\,\partial_\mu
\ln\Xi$ 再求一次导。(b) $\langle N\rangle \propto
\eu^{\beta\mu}$：$\partial_\mu\langle N\rangle = \beta\langle
N\rangle$，故 $\operatorname{Var}N = \langle N\rangle$。(c) 该题中的
方差为 $\bar n$ 的高斯分布，正是这一泊松陈述
在子体积中所见的情形。(d) $\operatorname{Var}n =
\bar n(1 - \bar n)$：该能级必定为空或必定
被占据时为零——泡利阻塞在费米海深处抑制了涨落。
\end{solution}

\begin{solution}{exo:b3:grand-canonical:8}
(a) 溶质稀释水的熵，从而降低其 $\mu$；压强
提高 $\mu$：二者平衡给出 $\Pi = ck_{\text{B}}T$。(b) $\Pi =
\num{3e26} \times \num{1.38e-23} \times 310 \approx
\qty{1.3}{MPa}$——相当于你血管内的十三个大气压，
并且经常保持平衡。(c) 纯水的 $\mu$ 高于细胞
内部：水涌入，细胞破裂（溶血）——
因此输液应为生理盐水，绝不用纯水。(d) 从根到叶的
$\mu$ 梯度（溶质、叶片上的蒸发）把水向上拉，
以渗透压和毛细压强差的形式，可达数十米，
且无需任何运动部件。
\end{solution}

\begin{solution}{exo:b3:grand-canonical:9}
(a) $\mu$ 处处相等：$k_{\text{B}}T\ln n(h) + mgh = $ 常数，即\
气压指数律。(b) $n(r) \propto
\eu^{+m\omega^2r^2/2k_{\text{B}}T}$：较重的组分聚集在
边缘。(c) $\exp[\Delta m\,\omega^2r^2/2k_{\text{B}}T] =
\exp[\num{5e-27} \times \num{3.6e5}/(2 \times \num{4.4e-21})]
\approx \eu^{0.2} \approx 1.2$ 每级。(d) 从天然
$0.7\%$ 提浓到反应堆级或武器级，需要多个
$1.2$ 的连乘因子：因此需要数千台离心机
组成的级联——这也是为何级联计数属于核外交范畴。
\end{solution}

\begin{solution}{exo:b3:grand-canonical:10}
(a) 同正文：对 $\mu$ 的平衡式取指数，并代入理想气体化学势
及原子的 $-E_{\text{I}}$ 偏移量。(b)
令 $n_{\text{e}} = n_{\text{p}} = xn$，$n_{\text{H}} =
(1-x)n$：得到所给的商。(c) 在 \qty{5800}{K} 下：
$\lambda_{T,\text{e}} \approx \qty{1.0}{nm}$，
$\eu^{-27.2} = \num{1.5e-12}$，$1/n\lambda^3 \approx \num{1.1e4}$：
$x^2/(1-x) \approx \num{1.6e-8}$，$x \approx \num{1.3e-4}$——
太阳光球层是 $99.99\%$ 的中性氢。(d) 在
\qty{10000}{K}（$n$ 不变）下：$x \approx 6\%$——
温度仅增加 $1.7$ 倍，电离度即升高两个半数量级；
在光球层较低的电子密度下，这种上升还会更早出现。萨哈
转变是陡崖式的。
\end{solution}

\begin{solution}{exo:b3:grand-canonical:11}
(a) $\num{1.7e-7}$：比阈值低七个数量级——空气
完全可按经典处理。(b) 令 $n\lambda_T^3 = 1$：$T_{\text{deg}} =
h^2n^{2/3}/2\pi mk_{\text{B}} \approx \qty{6e4}{K}$，即铜中电子的
简并温度——你摸过的每一种金属都携带
一种深度量子的电子气。(c) 铷在 \qty{e20}{m^{-3}} 时：
$T_{\text{deg}} \sim \qty{e-7}{K}$——正是
\cref{ch:b3:quantum-statistics}中的百纳开尔文前沿。(d) 金属中的电子：
始终是量子的；实验室原子：仅在纳开尔文量级；
空气：从不——简并程度按 $n^{2/3}/m$ 变化。
\end{solution}

\begin{solution}{exo:b3:grand-canonical:12}
(a) 当 $\mu = 0$ 时，电子对的数密度为受玻尔兹曼压低的
$\lambda_T^{-3}\eu^{-\beta m_{\text{e}}c^2}$，前提是电子对已非相对论化。
(b) $m_{\text{e}}c^2/k_{\text{B}} =
\qty{5.9e9}{K}$，约在宇宙起始后一秒达到。
(c) 从 $\qty{6e9}{K}$ 冷却至 $\qty{2e9}{K}$，指数乘以三：
仅指数项就使密度由 $\eu^{-2} \to$
$\eu^{-6}$ 下降，计入前置因子后再降数个数量级——正电子
在数秒之内便消失殆尽，几乎所有电子
都与之湮灭，只余十亿分之一。(d) 湮灭光子
汇入辐射浴，使其略微升温，其拉伸后的残余
即为宇宙微波背景
——这一事件的\qty{2.7}{K}“回执”。
\end{solution}

\begin{solution}{pb:b3:grand-canonical:1}
\textbf{1.} $n_2/n_1 = 4\,\eu^{-10.2\,\text{eV}/k_{\text{B}}T}$。
\textbf{2.} $\num{2e-10}$、$\num{2.9e-5}$、$\num{1.1e-2}$：急剧
上升——温度越高越有利于激发 $n = 2$ 能级。
\textbf{3.} 萨哈方程使氢电离殆尽：高于
$\sim\qty{10000}{K}$ 时\emph{中性}比例骤降，
巴耳末线也随之消失。
\textbf{4.} 一个因子随 $T$ 上升，另一个下降：二者
乘积必在两种极限之间取峰值。
\textbf{5.} 令 $\lambda_{T,\text{e}}^{-3}\eu^{-E_{\text{I}}/
k_{\text{B}}T} = n_{\text{e}} = \qty{e20}{m^{-3}}$，解得 $T
\approx \qty{9600}{K}$。
\textbf{6.} 接近 \qty{10000}{K} 时——恰是 A 型星，
其光谱正是氢线的典型范例。
\textbf{7.} 在 \qty{3000}{K} 时，没有任何粒子能达到 $n = 2$
（$\eu^{-39}$），而脆弱的分子如 TiO 得以存续，并以带状谱吸收：
冷星是分子星，氢几乎不显现。
\textbf{8.} 氢被电离至不可见；氦的
\qty{24.6}{eV} 能级梯，此时才具备 $k_{\text{B}}T$ 与萨哈
条件，以布居并保持其吸收态。
\textbf{9.} 钙的 \qty{6.1}{eV} 电离能远低于
氢的 \qty{13.6}{eV}：在 \qty{5800}{K} 时萨哈方程使钙几乎
完全电离（此时 Ca$^+$ 谱线从其基态吸收——无玻尔兹曼代价），
而中性氢停留在 $n = 1$，
对巴耳末线不可见。
\textbf{10.} 太阳主要由钙和铁构成，与地壳相似——这正是
佩恩所继承的、占统治地位的假设。
\textbf{11.} 将各谱线强度除以其可见度因子后，
剧烈的差异消失了：大多数元素呈现相近的
丰度比例——只有氢（以及氦）例外，
其丰度竟高出约\emph{一百万}倍。
\textbf{12.} 当时人们假定宇宙化学与地球岩石相似；
罗素敦促她把氢的结果称为“几乎
肯定不真实”——然而他四年后自己发表了同样的结论，
并注明了她的功劳。佩恩是对的。
\textbf{13.} $4\eu^{-10.2/0.827} = 4\eu^{-12.3} \approx
\num{1.8e-5}$。
\textbf{14.} 乘积 $\approx \num{9e-6}$、$\num{1.5e-6}$、
$\num{3.5e-7}$：最大值位于 $\approx\qty{9600}{K}$——即
A 型星的峰值，由两个指数式再现。
\textbf{15.} 在高温一侧，萨哈指数直接湮灭了
中性粒子的布居；在低温一侧，玻尔兹曼因子只是
渐渐衰减：上方是悬崖，下方是斜坡。
\textbf{16.} 更高的 $n_{\text{e}}$ 促使复合（可供捕获的
伙伴更多）：半电离温度向\emph{更高}温度方向移动——
使反应更偏向束缚一侧，即勒夏特列原理
在 $\mu$ 语言中的体现。
\textbf{17.} 表面重力：致密矮星（高压）与臃肿的
巨星（低压）在同一温度下分离——光度等级
由谱线物理读出。
\textbf{18.} 等先验态 $\to$ 玻尔兹曼布居
$\to$ 萨哈电离 $\to$ 谱线强度比反推
$T$：一台没有活动部件、可在任何距离工作的温度计。
\textbf{19.} 温度——一个维度，从 M 型星的 \qty{3000}{K}
到 O 型星的 \qty{40000}{K}，取代了整套字母分类。
\textbf{20.} 可见度即布居，而布居是 $T$ 的指数函数：谱线
强度首先反映恒星的热力学性质，
其化学组成只有
经修正后才能测得。
\textbf{21.} 氢聚变作为恒星能源，以及星系的 21 厘米制图，
二者皆立足于佩恩的
丰度结果。
\textbf{22.} 她借助两种系综公式，将四十年的光谱分类
转化为温度测量，并重新衡量宇宙的组成——
这一切都出自一篇博士论文。
\textbf{23.} 宇宙复合：宇宙自身的
半电离发生于 $\sim\qty{3000}{K}$，其释放的光，
经千倍拉伸后，即 \qty{2.7}{K} 的微波
背景。
\textbf{24.} 电子（高温，\qty{e4}{K} 量级）与气体
及管壁（低温）交换能量太慢、难以达到平衡：两种温度
共存于同一管中——这是一种非平衡稳态，
这正是灯管发光而不熔化的原因。
\textbf{25.} 一个条件 $\mu_{\text{H}} = \mu_{\text{p}} +
\mu_{\text{e}}$，仅凭温度便将 OBAFGKM 排序，
使巴耳末峰值落在约 \qty{9600}{K} 的 A 型星上，并且——
在对强度作可见度修正后——揭示出一个
以氢为主的宇宙。
\end{solution}
